\begin{proposition}[the three logarithmic relative direct images]
\label{prop:log-relative-direct-images}
\Verified\ One has
\begin{align*}
 f_*\mathcal E&\simeq\mathcal O\oplus\mathcal O(-1),\\
 R^1f_*\mathcal E&\simeq
       \mathcal O^{\oplus2}\oplus\mathcal O(-1)^{\oplus2},\\
 R^2f_*\mathcal E&\simeq\mathcal O\oplus\mathcal O(-1).
\end{align*}
All three sheaves are locally free and commute with base change.
\end{proposition}

\begin{proof}
At the cusp, the toric coordinates give
\[
 \mathcal E|_{N_0}
 =\mathcal O_{N_0}[d\log x_1]\oplus
  \mathcal O_{N_0}[d\log x_2].
\tag{11.3}
\]
Indeed a quotient shear changes each $d\log x_i$ only by a holomorphic
multiple of $d\log t$, which is zero in $\mathcal E$.  The already proved
structure-sheaf calculation therefore gives free local direct images of
ranks $2,4,2$ at $p_0$.

We now keep the ramification at $p_j$.  Write
\[
 A=\mathbb C\{t\},\qquad C=\mathbb C\{s\},\qquad t=s^m,
 \qquad (m,a)=(3,2)\text{ or }(4,1).
\]
Cartan's analytic quotient theorem identifies local holomorphic functions on
the quotient with invariants \cite[\S1, Theorem~1, pp.~1--2]{Cartan1954Quotient}.
The action on the total space is free, so the quotient map is finite \'etale;
the corresponding locally free sheaves and their equivariant structures
descend by finite \'etale descent.  This sheaf statement is not being
attributed to Cartan's theorem on invariant functions.  Finite pushforward
is exact, and Reynolds averaging
$v\mapsto |G_j|^{-1}\sum_{g\in G_j}gv$ makes the invariant functor exact in
characteristic zero.  Hence
\[
 R^qf_*\mathcal E|_{D_j}
 =\bigl(\pi_*R^q\widetilde f_*\widetilde{\mathcal E}\bigr)^{G_j}.
\tag{11.3a}
\]
Choose equivariant frames $e_0,e_1$ of the relative cotangent and
$b_0,b_1$ of $R^1\widetilde f_*\mathcal O$ whose semilinear characters are
respectively
\[
 (1,\zeta^a),\qquad (1,\zeta^{-a});
\]
$b_0$ is the nowhere-zero class $\gamma^{0,1}$.  These may be chosen as
holomorphic constant-character frames, rather than merely as eigenvectors on
the central fibre.  Indeed, if $w$ lifts a $\nu$-eigenvector at $s=0$, then
\[
 P_\nu(w)=\frac1m\sum_{k=0}^{m-1}\nu^{-k}(g^*)^kw
\]
satisfies $g^*P_\nu(w)=\nu P_\nu(w)$ and has the same nonzero central value;
after shrinking the disc, the projected eigenvectors form a frame.  On the
marked torus family the projection formula gives
\[
 R^q\widetilde f_*\widetilde{\mathcal E}
 \simeq
 (\mathcal O e_0\oplus\mathcal O e_1)
 \otimes
 \bigwedge^q(\mathcal O b_0\oplus\mathcal O b_1),
 \qquad 0\le q\le2.
\]
Solving
$\nu\zeta^r=1$ for the invariant exponent gives the following \emph{actual}
$A$-bases:
\begin{align*}
 f_*\mathcal E:\quad& e_0,\ s^{m-a}e_1,\\
 R^1f_*\mathcal E:\quad&
 e_0b_0,\ s^ae_0b_1,\ s^{m-a}e_1b_0,\ e_1b_1,\\
 R^2f_*\mathcal E:\quad&
 s^ae_0(b_0\wedge b_1),\ e_1(b_0\wedge b_1).
\tag{11.4}
\end{align*}
For comparison, the variable tangent lattice is generated by $s^av_1$.
Thus evaluation of the descended variable cotangent on the descended
variable tangent is $s^m=t$, not a unit.  This proves explicitly why
\[
 f_*\mathcal E\longrightarrow(f_*T_{X/B})^\vee
\]
is multiplication by $\operatorname{diag}(1,t)$ locally, although
$\mathcal E^\vee\simeq T_{X/B}$ before direct image.  This last identity is
an identity of locally free relative sheaves, not an assertion that formation
of the kernel commutes with taking a tangent-space fibre.  For example, at a
triple chart
\[
 R=\mathbb C\{x,y,z\},\qquad t=xyz,
\]
one has
\[
 T_{X/B}=\ker\!\left[R^3\xrightarrow{(yz,xz,xy)}R\right].
\]
If $yz\,u+xz\,v+xy\,w=0$, reduction modulo $x$ and $y$ gives
$u=xU$ and $v=yV$, and cancellation in the domain gives
$w=-z(U+V)$.  Hence
\[
 T_{X/B}
 =R(x\partial_x-z\partial_z)\oplus
  R(y\partial_y-z\partial_z)
 =\mathcal E^\vee.
\]
The kernel of the pointwise linear map $df_x$ at the triple point is
three-dimensional, whereas the fibre of this sheaf kernel is
two-dimensional: tensoring the kernel sequence with the residue field is
not left exact because the image ideal $(xy,xz,yz)$ is not flat.  Thus a
section can be nowhere zero in $T_{X/B}=\mathcal E^\vee$ even when its image
in the ambient tangent-space fibre vanishes.

The invariant vertical field $\xi$ gives a trivial subline of
$V:=f_*T_{X/B}$; the quotient has the reciprocal automorphy of the Hodge
line $\mathfrak K\simeq\mathcal O(1)$ from
\cref{lem:hodge-line-audit}.  Since $H^1(\mathcal O(1))=0$, the resulting
extension splits:
\[
 V\simeq\mathcal O\xi\oplus\mathcal O(-1),
 \qquad V^\vee\simeq\mathcal O\oplus\mathcal O(1).
\]
Evaluation gives a generically invertible map
$f_*\mathcal E\to V^\vee$.  Formula (11.4) shows that its elementary
divisors are $(1,t)$ at each of $p_1,p_2$ and $(1,1)$ elsewhere.  Thus its
cokernel has total length two and
\(\deg(f_*\mathcal E)=1-2=-1\).  Evaluation on $\xi$ is stalkwise
surjective: the covector dual to the fixed eigenvector is a unit generator
at the cusp and in both finite invariant bases.  Its kernel is consequently
a line bundle of degree $-1$, and
\[
 0\longrightarrow\mathcal O(-1)\longrightarrow f_*\mathcal E
 \xrightarrow{\operatorname{ev}_\xi}\mathcal O\longrightarrow0.
\]
Because $H^1(\mathcal O(-1))=0$, this extension also splits, proving
\[
 f_*\mathcal E\simeq\mathcal O\oplus\mathcal O(-1).
\]
Also, by 
\cref{thm:structure-direct-images-audit},
\[
 G:=R^1f_*\mathcal O_X=\mathcal O\oplus\mathcal O(-1),
 \qquad R^2f_*\mathcal O_X=\det G=\mathcal O(-1).
\tag{11.5}
\]
Here is the global filtration implicit in those bases.  Choose the
degree-zero summand in the preceding split sequence so that its generator
evaluates to one on $\xi$; locally this is the fixed eigen-covector $e_0$.
Together with the fixed class $b_0=\gamma^{0,1}$, take the corresponding
saturated sublines in $f_*\mathcal E$ and $G$.  Formula (11.3) at the cusp
and the bases (11.4) at the finite points show that these sublines are locally
generated by $e_0,b_0$ and that their quotients are locally generated by
$s^{m-a}e_1,s^ab_1$; hence both inclusions and both quotients are locally
free.  The evaluation splitting and Hodge-line calculation identify the two
filtrations as
\[
 0\subset\mathcal O\subset f_*\mathcal E,\qquad
 0\subset\mathcal O\subset G,\qquad
 \operatorname{gr}(f_*\mathcal E)=
 \operatorname{gr}(G)=\mathcal O\oplus\mathcal O(-1).
\]
Over the smooth locus the cup-product identification is
$R^1f_*\mathcal E=(f_*\mathcal E)\otimes G$.  Formula (11.4) extends its four
character lines as saturated rank-one subsheaves over all of $B$.  Three
associated-graded lines are generated locally by
$e_0b_0,s^ae_0b_1,s^{m-a}e_1b_0$ and have degrees $0,-1,-1$.  On the fourth,
the separately descended factors would
give
\[
 (s^{m-a}e_1)(s^ab_1)=t(e_1b_1),
\]
whereas $e_1b_1$ itself is invariant.  Thus that graded line is the upper
modification of $\mathcal O(-2)$ at $p_1+p_2$, hence is $\mathcal O$.
Similarly, the two graded lines of $R^2f_*\mathcal E$ are
$\mathcal O(-1)$ and
$\mathcal O(-2)(p_1+p_2)=\mathcal O$.  Every extension between the resulting
degree-zero and degree-minus-one lines splits: its $\operatorname{Ext}^1$
is one of $H^1(\mathcal O(-1))$, $H^1(\mathcal O)$, or
$H^1(\mathcal O(1))$, all zero.  This proves the three asserted splittings.
This is also the elementary filtered form of the
Birkhoff--Grothendieck classification
\cite[Theorem~2.1, p.~126]{Grothendieck1957}; the divisor modifications and
$\operatorname{Ext}^1$ identifications may equivalently be read in
\cite[Tags 0C6K, 0B3A, and 01XT]{StacksProject}.

Finally, (11.3a)--(11.4) make every direct image free of ranks $2,4,2$ at
the finite points; (11.3) does the same at the cusp, and the smooth locus is
standard.  The multiple-fibre base change can also be read directly from the
invariant modules.  The Reynolds idempotent is $A$-linear, so invariants
commute with arbitrary finite $A$-base change.  For the scheme-theoretic
special fibre,
\[
 C\otimes_AA/(t)=C/(s^m),
 \qquad
 \bigl((C/(s^m))v\bigr)^G=\mathbb C\,s^rv
\]
on a character line whose invariant exponent satisfies $0\le r<m$.
Every basis vector in (11.4) therefore survives reduction modulo $t$, and
the special-fibre dimensions are exactly $2,4,2$.  No $R^q$ with $q>2$
occurs.  Descending induction in analytic cohomology and base change then
makes every base-change map an isomorphism, including for the nonreduced
finite fibres; precisely the same criterion, with citation, is recorded in
the proof of
\cref{thm:structure-direct-images-audit}; see in particular
\cite[Chapter~III, \S3, Theorem~3.4 and Corollary~3.10]{BanicaStanasila1976}.
\end{proof}
